Un pequeño equipo de ensayo idealizado tiene temperatura uniforme \(T\), capacidad térmica constante \(C=600\,\mathrm{J/K}\) y está en aire a \(20\,{}^\circ\mathrm C\). Elevar su temperatura en \(\Delta T\) requiere una energía \(C\Delta T\). La potencia que cede al aire es \(H(T-20)\), con \(H=10\,\mathrm{W/K}\); las diferencias de temperatura se expresan en kelvin o, con el mismo valor numérico, en grados Celsius.

Un calefactor aporta \(400\,\mathrm W\) cuando está encendido y cero cuando está apagado. Un termostato mide directamente la temperatura del equipo: apaga el calefactor al llegar a \(40\,{}^\circ\mathrm C\), lo enciende al bajar a \(30\,{}^\circ\mathrm C\) y mantiene la posición anterior del interruptor entre ambos umbrales. La conmutación no tiene retardo ni aporta energía de forma instantánea. No hay otras fuentes ni pérdidas de calor. Los datos son sintéticos.

En \(t=0\), el equipo está a \(30\,{}^\circ\mathrm C\) y el calefactor acaba de encenderse. Elija la información necesaria para describir su estado y formule el modelo. Obtenga la temperatura durante un ciclo completo, su duración y la fracción de tiempo en que el calefactor permanece encendido. Calcule también la temperatura media de un ciclo mediante el balance de energía, sin integrar por separado las dos expresiones de temperatura. Explique si conocer únicamente \(T=35\,{}^\circ\mathrm C\) permite determinar la derivada temporal de la temperatura.
